\documentclass{ohm_project_description}
\usepackage[english]{babel}

\projecttitle{M-APR: Handling of Flexible Materials Using Robot Arms}
\projectauthor{TTZ Nürnberger Land / TH Nürnberg}
\projectdate{2026}

\begin{document}

\maketitle
\thispagestyle{fancy}

Flexible materials such as films, textiles, plastics, cables or limp components are gaining increasing importance across many industries. With rising production automation, reliable and precise handling of these materials is becoming a key challenge, as conventional robot systems are primarily designed for rigid workpieces.

This M-APR project (3 semesters, starting WiSe 2026/2027) addresses the handling of flexible materials by robot arms. The exact focus is agreed jointly and can cover \textbf{hardware aspects} (specialised gripping technologies) or \textbf{software/AI aspects} (simulation, learning-based control, Deep Reinforcement Learning, Vision-Language-Action models). The work concludes with implementation and evaluation in a physical robot cell at the TTZ Nürnberger Land.

\section*{Work Packages (selection by focus)}
\begin{itemize}[leftmargin=0.5cm]
    \setlength\itemsep{.1em}
    \item Analysis of gripping principles for flexible materials (vacuum, adhesive, form-adaptive)
    \item Development and fabrication of a specialised gripper (hardware focus)
    \item Building a simulation environment for deformable objects (software focus)
    \item Model-based control or Deep Reinforcement Learning (DRL)
    \item Vision-Language-Action (VLA) models for flexible manipulation
    \item Sim-to-real transfer and evaluation in the robot cell at TTZ
\end{itemize}

\section*{Requirements}
\begin{itemize}[leftmargin=0.5cm]
    \setlength\itemsep{.1em}
    \item Bachelor's degree in electrical engineering, computer science, mechatronics, robotics or equivalent
    \item Programming in C++ and/or Python
    \item Basic knowledge of machine learning and robotics
    \item Independent and goal-oriented working style
\end{itemize}

\vspace{0.5cm}
This topic is designed as an \textbf{M-APR (3 semesters)}, but can also be completed as a project or master's thesis subject to agreement.

\vfill
\textcolor{ohm_red}{\rule{\linewidth}{0.4mm}}
\textbf{\textcolor{ohm_red}{TTZ Nürnberger Land / Mobile Robotics Lab}} \\
\textbf{Location:} Martin-Luther-Str.\ 18, 91207 Lauf an der Pegnitz \\[4pt]
\begin{tabular}{@{}ll}
\textbf{Supervisor:}    & Prof. Dr. Christian Pfitzner \\
\textbf{E-Mail:}        & \href{mailto:christian.pfitzner@th-nuernberg.de}{christian.pfitzner@th-nuernberg.de} \\
\textbf{Co-Supervisor:} & M.\,Sc.\ Patrick Fußy \\
\textbf{E-Mail:}        & \href{mailto:patrick.fussy@th-nuernberg.de}{patrick.fussy@th-nuernberg.de} \\
\end{tabular}

\end{document}
